\documentclass[12pt]{article}
\usepackage[utf8]{inputenc}
\usepackage[T1]{fontenc}
\usepackage{hyperref}
\title{Zo'\'e Language Resources for K-12 Curriculum Development}
\author{Compiled by Claude Sonnet 5.5 (claude-sonnet-5-5) for Andr\'as Kornai}
\date{}
\begin{document}
\maketitle

\section{Introduction}
Zo'\'e (ISO 639-3: pto; Glottocode zoee1240) is a Tupi-Guarani language of the Tupian family, spoken in northern Par\'a, Brazil. Also called Zo\'e, Bur\'e, Poturu and Poturujara. Sources give about 150 speakers (1998, Wikipedia) and a people of about 350 (Iep\'e, 2025). The community is reported as largely monolingual in Zo'\'e, with a few young people having picked up some Portuguese. FUNAI restricts access to the Zo'\'e territory, so outside research and materials are limited.

\section{Primary Educational Resources}
No primers, textbooks or school programmes in Zo'\'e were found in the searches carried out. The language is not known to be taught in schools.

\section{Linguistic and Reference Materials}
\begin{itemize}
\item Wikipedia article on the Zo'\'e language: \url{https://en.wikipedia.org/wiki/Zo%CA%BC%C3%A9_language}
\item Endangered Languages Project entry for Zo'\'e: \url{https://www.endangeredlanguages.com/lang/1632?hl=en}
\end{itemize}
No dictionary or grammar was identified in these search results.

\section{Religious and Community Resources}
\begin{itemize}
\item Povos Ind\'igenas no Brasil (Instituto Socioambiental) page on the Zo'\'e people: \url{https://pib.socioambiental.org/en/Povo:Zo'é}
\item Wikipedia article on the Zo'\'e people: \url{https://en.wikipedia.org/wiki/Zo%CA%BC%C3%A9}
\end{itemize}

\section{Digital and Audio Resources}
No audio, video, corpus or software resources specific to Zo'\'e were found in the searches.

\section{Curriculum Development Context}
The language is the first language of the whole community, which makes it a natural medium of instruction, but access is controlled by FUNAI and the community has only recently left isolation. Any curriculum work would need to proceed through the community and the relevant Brazilian authorities. The sources found offer background only, not teaching materials.

\end{document}